\chapter{Abdominoplasty}

\section*{What Is This Surgery?}
Abdominoplasty, commonly known as a tummy tuck, is a major cosmetic surgical procedure designed to remove excess skin and fat from the middle and lower abdomen and to tighten the underlying abdominal wall musculature. The operation aims to create a flatter, firmer abdominal contour, particularly when diet, exercise, and lifestyle modifications have failed to address significant skin laxity, subcutaneous fat deposits, and muscular diastasis. This procedure predominantly targets the anterior abdominal wall, extending from the costal margins and xiphoid process superiorly to the pubic region inferiorly. It is frequently performed following substantial weight loss, after multiple pregnancies, or as a key component of a comprehensive body-contouring plan. Beyond its aesthetic benefits, abdominoplasty can alleviate functional symptoms such as chronic intertrigo, skin irritation, and difficulties with hygiene caused by a large, overhanging pannus, underscoring its importance as more than a purely cosmetic intervention.

\section*{Alternative Names}
\begin{itemize}
  \item Tummy tuck
  \item Abdominal lipectomy
  \item Abdominal dermolipectomy
  \item Rectus plication with abdominoplasty
  \item Panniculectomy (a related but more limited procedure focused on functional excision)
\end{itemize}

\section*{Relevant Surgical Anatomy}
A thorough understanding of the anterior abdominal wall anatomy is paramount for safe and effective abdominoplasty. The layers encountered, from superficial to deep, include the skin, subcutaneous fat (Camper's fascia and Scarpa's fascia), the superficial muscular aponeurotic system, and the rectus abdominis muscles encased within their fascial sheaths. The rectus abdominis muscles run vertically from the xiphoid process and costal cartilages to the pubic symphysis, separated in the midline by the linea alba. Diastasis recti, a common indication for abdominoplasty, is a widening of the linea alba, not a true hernia. The umbilicus is a critical landmark, typically located at the level of the L3-L4 vertebrae, and its stalk connects to the deep fascia.

The vascular supply to the abdominal skin flap is crucial. The primary blood supply to the elevated flap comes from the superior epigastric vessels (branches of the internal thoracic artery) and the deep inferior epigastric vessels (branches of the external iliac artery), which run within the rectus sheath and provide perforators to the overlying skin. The superficial epigastric and superficial circumflex iliac arteries, originating from the femoral artery, contribute to the superficial vascular network but are often transected during flap elevation. Venous drainage largely parallels the arterial supply. Sensory innervation to the abdominal wall is provided by the anterior cutaneous branches of the intercostal nerves (T7-T11) and the iliohypogastric and ilioinguinal nerves (T12-L1), which are often disrupted during flap elevation, leading to temporary or permanent numbness. Lymphatic drainage follows the superficial venous system, primarily to the inguinal lymph nodes. Surgical landmarks include the anterior superior iliac spines, the pubic symphysis, and the costal margins, which define the boundaries of the surgical field.

\section*{Surgical Principle \& Rationale}
The fundamental operative concept of abdominoplasty revolves around the comprehensive reshaping of the anterior abdominal wall by excising redundant skin and subcutaneous fat, while simultaneously restoring the structural integrity and muscular support of the abdominal core. The surgeon meticulously elevates a large skin and fat flap from the underlying aponeurotic layer, typically up to the costal margins. This allows for direct access to the rectus abdominis muscles and their fascial sheaths. The core mechanism of correction involves the plication (suturing together) of the stretched rectus abdominis muscles and their investing fascia in the midline, effectively closing the diastasis recti that commonly contributes to a protruding abdomen and weakened core. This muscle tightening not only narrows the waistline but also recreates a robust internal support sling, improving core stability and potentially alleviating back pain.

The key technical goals of abdominoplasty are multifaceted: to achieve a smooth, tension-free closure with a low transverse scar that can be concealed by underwear; to meticulously preserve the blood supply to the elevated skin flap to prevent necrosis; and to accurately reposition the umbilicus through a new opening in the redraped skin, ensuring a natural and aesthetically pleasing appearance. The rationale extends beyond mere aesthetics, aiming to improve functional aspects such as posture, core strength, and relief from skin-fold related issues. The procedure is a balance of maximal tissue removal, optimal muscle tightening, and careful preservation of vascularity to achieve a durable and aesthetically superior result.

\section*{History \& Surgical Evolution}
The modern abdominoplasty has a rich history, evolving significantly over the past century through a series of surgical innovations. Early attempts at abdominal lipectomy date back to the late 19th century, with pioneers like Demars and Marx in France (1890) and Kelly in the United States (1899) describing simple elliptical excisions of large abdominal pannuses, primarily for functional relief rather than aesthetic contouring. These initial procedures were often limited in scope and associated with significant complications due to a lack of understanding of flap vascularity.

A major turning point occurred in the mid-20th century. In the late 1950s and 1960s, surgeons began to understand the importance of more extensive skin flap undermining and introduced the concept of umbilical transposition, recognizing the need to preserve the navel for a natural appearance. Dr. Ivo Pitanguy, a Brazilian plastic surgeon, is widely credited with refining and popularizing the standard low transverse incision with muscle plication in the 1960s and 1970s. His systematic approach established many of the principles that remain fundamental to abdominoplasty today, emphasizing both aesthetic contour and functional restoration. Subsequent advances included the integration of liposuction-assisted techniques in the 1980s, which allowed for improved contouring of the flanks and upper abdomen, and the development of variations such as the fleur-de-lis abdominoplasty for patients with significant vertical skin excess, particularly after massive weight loss. These continuous refinements have progressively reduced complication rates while significantly enhancing the cosmetic and functional outcomes of the procedure.

\section*{Clinical Indications}
Abdominoplasty is indicated for individuals presenting with specific anatomical and functional concerns related to the anterior abdominal wall. Key clinical indications include:

\begin{itemize}
  \item Significant excess abdominal skin and laxity, often observed after massive weight loss (e.g., post-bariatric surgery) or multiple pregnancies, which is unresponsive to diet and exercise.
  \item Diastasis recti, characterized by a separation of the rectus abdominis muscles in the midline, leading to abdominal wall weakness, a protruding belly, and sometimes associated back pain.
  \item Chronic skin irritation, intertrigo, rashes, or fungal infections occurring beneath an overhanging abdominal pannus, which can be debilitating and difficult to manage with conservative measures.
  \item Persistent abdominal bulge and skin redundancy that causes discomfort, interferes with clothing fit, or negatively impacts body image, despite achieving a stable and healthy weight.
  \item Functional symptoms such as lower back strain or postural issues attributed to the weight and laxity of a heavy abdominal pannus.
  \item Patients seeking improved abdominal contour and body image who have realistic expectations and are committed to maintaining a stable weight postoperatively.
  \item In some cases, to facilitate personal hygiene and mobility when an extensive pannus creates physical impediments.
\end{itemize}

\section*{Clinical Positioning}
Abdominoplasty is primarily positioned as an elective body-contouring procedure, typically performed in patients who have achieved and maintained a stable weight for several months and are in good general health. It is generally considered a primary procedure when addressing aesthetic concerns or functional issues like chronic intertrigo in otherwise healthy individuals. However, its role often shifts to a secondary or reconstructive step, particularly following massive weight loss, such as after bariatric surgery. In this context, abdominoplasty completes the patient's physical transformation, addressing the residual skin laxity that weight loss alone cannot resolve.

While predominantly elective, the procedure can be considered functionally reconstructive when a large, symptomatic pannus causes significant physical discomfort, hygiene issues, or mobility limitations. It is rarely combined with other intra-abdominal surgeries due to the increased risk of wound contamination and infection, as well as the potential for prolonged operative time and recovery. Therefore, it is almost always performed as a standalone procedure, allowing for optimal focus on the abdominal contour and minimizing surgical risks. Patient selection is critical, ensuring that candidates are psychologically prepared, have realistic expectations, and understand the commitment required for postoperative care and long-term maintenance of results.

\section*{Surgical Technique \& Procedure}
The abdominoplasty procedure is performed under general anesthesia, with the patient positioned supine on the operating table. The patient is often placed in a semi-flexed position (jack-knife) during closure to reduce tension on the incision.

\begin{enumerate}
  \item  \textbf{Anesthesia and Positioning:} General endotracheal anesthesia is administered. The patient is positioned supine, often with the operating table capable of flexion to aid in closure. Preoperative markings are meticulously drawn to delineate the planned incision lines, areas of skin excision, and the new umbilical position.
  \item  \textbf{Incision:} A low transverse incision is made just above the pubic hairline, extending laterally towards the anterior superior iliac spines. The length and shape of this incision vary based on the amount of skin laxity. For patients with significant vertical skin excess, a vertical midline incision may also be incorporated (fleur-de-lis abdominoplasty). A circumumbilical incision is made to free the umbilicus from the surrounding skin.
  \item  \textbf{Flap Elevation:} The skin and subcutaneous fat (dermofat flap) are carefully elevated from the underlying rectus fascia. This dissection proceeds superiorly, typically to the level of the costal margins and xiphoid process, creating a large vascularized flap. Meticulous hemostasis is maintained throughout this stage using electrocautery to minimize bleeding and hematoma formation.
  \item  \textbf{Rectus Plication:} Once the flap is elevated, the rectus abdominis muscles and their fascial sheaths are exposed. If diastasis recti is present, the rectus muscles are plicated (sutured together) in the midline using strong, non-absorbable sutures. This step tightens the abdominal wall, narrows the waist, and restores core integrity. Multiple rows of sutures may be used for a robust repair.
  \item  \textbf{Umbilical Transposition:} The umbilicus, which remains attached to its stalk originating from the deep fascia, is brought through a new opening created in the redraped abdominal skin flap. Its position is carefully chosen to ensure a natural appearance, typically at the level of the anterior superior iliac spines.
  \item  \textbf{Skin Excision and Redraping:} The elevated skin flap is then pulled inferiorly and redraped over the tightened abdominal wall. The excess skin and fat in the lower abdominal region, below the new umbilical position, are precisely excised. The amount of tissue removed can be substantial, often weighing several kilograms.
  \item  \textbf{Closure:} The wound is closed in multiple layers. The deep fascial layer is approximated first, followed by the subcutaneous tissue, and finally the skin. This layered closure helps distribute tension and minimize scarring. One or two closed-suction drains (e.g., Jackson-Pratt drains) are typically placed beneath the flap to prevent seroma formation.
  \item  \textbf{Dressing and Compression:} A sterile dressing is applied, and a compression garment is placed immediately postoperatively to provide support, reduce swelling, and aid in contouring.
\end{enumerate}

\section*{Preoperative Workup \& Preparation}
Comprehensive preoperative workup and preparation are essential to optimize patient safety and surgical outcomes for abdominoplasty. This process begins with a thorough medical history and physical examination to identify any comorbidities or risk factors. Patients must demonstrate a stable weight for at least six months prior to surgery, as significant weight fluctuations can compromise the long-term result. Smoking cessation is an absolute requirement for a minimum of four to six weeks preoperatively, as smoking severely impairs skin flap vascularity and dramatically increases the risk of skin necrosis and wound healing complications.

\begin{itemize}
  \item \textbf{Medical Evaluation:} A complete medical evaluation, including cardiac and pulmonary clearance, is obtained, especially for patients with pre-existing conditions. This may involve an ECG, chest X-ray, and consultations with specialists.
  \item \textbf{Laboratory Studies:} Routine preoperative laboratory tests include a complete blood count (CBC), basic metabolic panel, coagulation profile (PT/INR, PTT), and sometimes a urinalysis.
  \item \textbf{Medication Adjustment:} Patients are instructed to discontinue anticoagulant and antiplatelet medications (e.g., aspirin, NSAIDs, warfarin, novel oral anticoagulants) typically 7-10 days before surgery to minimize bleeding risk. Herbal supplements with anticoagulant properties should also be stopped.
  \item \textbf{NPO Status:} Patients are required to be NPO (nil per os) for at least 6-8 hours prior to surgery to prevent aspiration during anesthesia.
  \item \textbf{Prophylactic Antibiotics:} Intravenous prophylactic antibiotics are administered within one hour of incision to reduce the risk of surgical site infection.
  \item \textbf{Informed Consent and Counseling:} Detailed informed consent is obtained, covering the procedure, potential risks, benefits, alternative treatments, and expected recovery. Realistic expectations regarding scarring, contour, and recovery timeline are thoroughly discussed. Psychological assessment for body dysmorphic disorder may be considered in selected cases.
\end{itemize}

\section*{Contraindications \& Precautions}
Careful patient selection is critical for abdominoplasty to ensure safety and satisfactory outcomes.

\textbf{Absolute Contraindications:}
\begin{itemize}
  \item Uncontrolled chronic medical conditions (e.g., severe cardiovascular disease, uncontrolled diabetes mellitus, significant pulmonary disease, active autoimmune disorders).
  \item Active smoking that cannot be ceased for the recommended preoperative period (typically 4-6 weeks).
  \item Body Mass Index (BMI) above a certain threshold (e.g., >30-35 kg/m\textasciicircum2), which significantly increases operative risks and complication rates.
  \item Unrealistic patient expectations or untreated psychological conditions such as body dysmorphic disorder.
  \item Plans for future pregnancy, as subsequent pregnancies would likely reverse the muscle tightening and compromise the aesthetic result.
  \item Active infection or malignancy.
\item Significant coagulopathy or bleeding disorders that cannot be safely managed.
\end{itemize}

\textbf{Relative Contraindications \& Precautions:}
\begin{itemize}
  \item Significant previous abdominal surgeries with extensive scarring, which may compromise the vascularity of the abdominal flap and increase the risk of skin necrosis.
  \item Morbid obesity (BMI >40 kg/m\textasciicircum2), where a panniculectomy might be a safer initial step, or further weight loss is recommended.
  \item Poor nutritional status.
  \item History of deep vein thrombosis (DVT) or pulmonary embolism (PE), requiring aggressive prophylactic measures.
  \item Certain medications that cannot be safely discontinued.
  \item Concurrent intra-abdominal pathology that might require a separate surgical approach or increase the risk of contamination.
  \item Patients with a history of keloid or hypertrophic scarring.
\item Unstable weight, with plans for further significant weight loss.
\end{itemize}

\section*{Complications \& Risks}
As with any major surgical procedure, abdominoplasty carries inherent risks and potential complications, which can be broadly categorized into general surgical risks and procedure-specific risks.

\textbf{General Surgical Complications:}
\begin{itemize}
  \item \textbf{Bleeding and Hematoma:} Accumulation of blood under the skin flap, potentially requiring surgical drainage.
  \item \textbf{Infection:} Surgical site infection, ranging from superficial cellulitis to deep abscess, requiring antibiotics or drainage.
  \item \textbf{Anesthesia Risks:} Adverse reactions to anesthesia, including nausea, vomiting, respiratory complications, or allergic reactions.
  \item \textbf{Deep Vein Thrombosis (DVT) and Pulmonary Embolism (PE):} Formation of blood clots in the legs that can travel to the lungs, a potentially life-threatening complication. Prophylaxis is crucial.
  \item \textbf{Pain:} Postoperative pain, which is managed with analgesics.
\end{itemize}

\textbf{Procedure-Specific Complications:}
\begin{itemize}
  \item \textbf{Seroma:} The most common complication, involving the accumulation of serous fluid under the skin flap, often requiring repeated aspirations or drain reinsertion.
  \item \textbf{Skin Flap Necrosis:} Partial or full loss of the elevated skin flap due to compromised blood supply, particularly in smokers or patients with extensive undermining. This can lead to delayed wound healing, open wounds, and significant scarring.
  \item \textbf{Delayed Wound Healing and Dehiscence:} Separation of wound edges, especially at the incision line, which may require prolonged wound care or revision.
  \item \textbf{Altered Sensation:} Numbness or hypersensitivity of the abdominal skin, which can be temporary or permanent, due to nerve transection or injury during flap elevation.
  \item \textbf{Unfavorable Scarring:} Wide, hypertrophic, or keloid scars, or an aesthetically displeasing umbilical scar. Scar revision may be necessary.
  \item \textbf{Asymmetry or Irregular Contour:} Unevenness or irregularities in the abdominal contour, potentially requiring revision surgery.
  \item \textbf{Umbilical Complications:} Necrosis of the umbilical stalk, stenosis of the umbilical opening, or an unnatural appearance of the umbilicus.
  \item \textbf{Fat Necrosis:} Hardened areas under the skin due to localized fat cell death, which can be palpable and sometimes painful.
  \item \textbf{Recurrence of Diastasis Recti:} Although rare, the rectus plication can fail, leading to a recurrence of abdominal wall laxity.
\end{itemize}

\section*{Postoperative Recovery Timeline}
The postoperative recovery from abdominoplasty is a gradual process that requires patience and adherence to surgical instructions. Immediately after surgery, patients are transferred to the Post-Anesthesia Care Unit (PACU) for close monitoring of vital signs, pain levels, and drain output.

During the first 24-48 hours, patients typically experience moderate to severe pain, managed with prescribed oral or intravenous analgesics. They are encouraged to ambulate gently with assistance to prevent DVT, maintaining a slightly flexed position to reduce tension on the incision. Drains are monitored for output, and the compression garment remains in place. Most patients are discharged within one to two nights, depending on pain control and drain output. For the first 1-2 weeks, activity is significantly restricted; heavy lifting, strenuous exercise, and activities that strain the abdominal muscles are strictly forbidden. Patients should continue to walk in a slightly bent position. Diet is advanced as tolerated, typically starting with clear liquids and progressing to a regular diet. Drain removal usually occurs when output falls below a specified threshold, often within 7-14 days. The compression garment is worn continuously for several weeks (typically 4-6 weeks) to minimize swelling and support the healing tissues. Long-term recovery involves a progressive return to normal activities and exercise over 6-8 weeks. Swelling can persist for several months, and the final contour and scar maturation may take up to a year or more to fully develop. Sensation in the abdominal skin often returns slowly and may not fully normalize.

\section*{Surgical Findings \& Pathology}
During abdominoplasty, the surgeon documents several key findings. These typically include the extent of skin and subcutaneous fat redundancy, the degree of diastasis recti (measured in centimeters), the presence of any umbilical hernias, and the quality of the abdominal wall fascia. The excised tissue, comprising excess skin and subcutaneous fat, is routinely sent for pathological examination.

The pathology report on the resected tissues usually describes benign adipose tissue and skin. The skin may show evidence of striae gravidarum (stretch marks), hyperpigmentation, or other benign dermatological changes. While rare, the pathologist will also screen for any incidental findings such as lipomas, sebaceous cysts, or other unexpected lesions within the excised tissue. The primary purpose of the pathological examination is to confirm the benign nature of the removed tissue and to rule out any unforeseen pathology, providing an important safety check for the patient. The surgeon's operative notes will detail the extent of muscle plication, the volume of tissue excised, and the successful transposition of the umbilicus, forming the basis for assessing the immediate surgical outcome.

\section*{Expected Postoperative Course}
A normal and successful postoperative course following abdominoplasty is characterized by a predictable progression of healing and resolution of initial symptoms. Patients can expect initial swelling and bruising across the abdomen, which gradually subsides over several weeks to months. Pain, initially moderate, should progressively decrease and be manageable with prescribed oral analgesics, typically resolving significantly within the first 1-2 weeks. The incision line should heal cleanly without signs of infection or dehiscence, and the drains, if placed, will be removed once output is minimal.

The abdomen will gradually flatten and firm up as swelling resolves and the muscle plication takes effect. While some numbness in the lower abdominal skin is common initially, sensation typically begins to return over several months, though it may not fully normalize. The surgical scar, initially red and raised, will soften and fade over 6-12 months, eventually becoming a fine, pale line. Patients should experience improved abdominal contour, a firmer core, and resolution of any pre-existing functional symptoms such as skin irritation or back pain. The expected outcome is a stable, aesthetically pleasing abdominal contour, provided the patient maintains a consistent weight and healthy lifestyle.

\section*{Postoperative Care \& Follow-up}
Diligent postoperative care and regular follow-up appointments are crucial for optimal healing and long-term success after abdominoplasty.

\begin{itemize}
  \item \textbf{Follow-up Visits:} Patients typically have their first postoperative visit within 5-10 days to assess wound healing, check for complications, and remove drains if output criteria are met. Subsequent visits are scheduled at 2-4 weeks, 3 months, 6 months, and 1 year to monitor progress.
  \item \textbf{Wound Care:} Instructions for incision care, including gentle cleaning and dressing changes, are provided. Patients are advised to keep the incision dry and clean.
  \item \textbf{Suture/Staple Removal:} If non-absorbable sutures or staples are used, they are typically removed at the first or second follow-up visit.
  \item \textbf{Compression Garment:} The compression garment must be worn continuously for 4-6 weeks, and sometimes longer, to reduce swelling, support the tissues, and aid in contouring.
  \item \textbf{Activity Restrictions:} Patients are advised to avoid strenuous activities, heavy lifting, and abdominal exercises for at least 6-8 weeks to protect the muscle plication and allow for proper healing. Gradual return to light activities is encouraged.
  \item \textbf{Scar Management:} Once the incision is fully healed, scar massage, silicone sheets, or topical scar creams may be recommended to optimize scar maturation and appearance.
  \item \textbf{Warning Signs:} Patients are educated on warning signs that necessitate immediate contact with the surgeon, including fever, increasing pain unresponsive to medication, excessive wound drainage, sudden swelling, redness spreading from the incision, or shortness of breath.
\end{itemize}
Long-term follow-up emphasizes maintaining a stable weight and a healthy lifestyle to preserve the surgical result and prevent recurrence of skin laxity or diastasis.

\section*{Epidemiology}
Abdominoplasty consistently ranks among the most frequently performed body-contouring procedures globally. Its incidence has seen a notable increase in recent decades, largely paralleling the rise in bariatric surgery and the subsequent demand for post-massive weight loss body contouring. The procedure is predominantly performed in adult women, often after pregnancy, which is a significant contributing factor to abdominal wall laxity and diastasis recti. However, it is also increasingly common in men who have experienced substantial weight loss.

The typical age range for patients undergoing abdominoplasty spans from the late 20s to the 60s, with peak incidence in mid-adulthood. While the procedure is associated with a low perioperative mortality rate in appropriately selected patients, morbidity is primarily dominated by wound complications such as seroma (reported in 10-30\% of cases), infection (1-5\%), and delayed wound healing. Serious complications like DVT/PE are rare but remain a concern, with rates typically below 1\%. The functional indications, such as relief from chronic intertrigo, also contribute significantly to the procedure's overall frequency, highlighting its dual role in both aesthetic and reconstructive surgery.

\section*{Outcomes \& Prognosis}
When performed in appropriately selected patients by experienced surgeons, abdominoplasty yields high rates of patient satisfaction and reliable, long-lasting improvements in abdominal contour and core strength. Studies consistently report patient satisfaction rates exceeding 90\%. The main determinants of a favorable outcome include a stable postoperative weight, complete cessation of smoking, good general health, and realistic patient expectations regarding scarring and the final aesthetic result.

Functionally, abdominoplasty effectively resolves symptoms such as chronic intertrigo, skin irritation, and difficulties with hygiene caused by an overhanging pannus. Many patients also report improvements in posture and a reduction in lower back pain due to the tightened abdominal musculature. The prognosis for maintaining the surgical result is excellent, provided the patient adheres to a healthy lifestyle and avoids significant weight fluctuations or future pregnancies. While the surgical scar is permanent, it typically matures and fades over time, becoming less conspicuous. Recurrence of diastasis recti is uncommon after a robust plication, but minor contour irregularities or scar revisions are occasionally required. Landmark studies and large cohort analyses consistently demonstrate the efficacy and safety of abdominoplasty as a transformative procedure for patients seeking both aesthetic enhancement and functional improvement of the abdominal wall.

\section*{References}
\begin{enumerate}
  \item American Society of Plastic Surgeons. Tummy Tuck Surgery Guide. 2024. Available from: \url{https://www.plasticsurgery.org/cosmetic-procedures/tummy-tuck}.
  \item Mathes SJ, Hentz VR, editors. Plastic Surgery. 2nd ed. Philadelphia, PA: Saunders Elsevier; 2006. Chapter 109: Abdominoplasty.
  \item Pitanguy I. Abdominal lipectomy: an approach to it through an analysis of 300 consecutive cases. Plast Reconstr Surg. 1967;40(4):384-391.
  \item MedlinePlus. Abdominoplasty (Tummy Tuck). U.S. National Library of Medicine; 2025. Available from: \url{https://medlineplus.gov/ency/article/002985.htm}.
\end{enumerate}

\clearpage